% !PRINTED 378
\weberpage{378}
%<<CONT
Wurzeln von (9) und nur diese, wenn man in (8) die Wurzeln
$x'$, $x''$ von (7) paarweise kombiniert, und daraus folgt (6).

Nach (6) ist nur noch nötig, die Funktion $\psi(D,m)$ für den
Fall zu bestimmen, daß $m=q^k$ eine Potenz einer in $Q$ nicht
aufgehenden Primzahl $q$ ist. Wir brauchen das Symbol $(D,n)$ in
dem in \S\,85 erklärten Sinn und unterscheiden drei Fälle:

1) Wenn $q$ in $\Delta$ aufgeht, so ist $(D,q)=0$, und die Kongruenz
 (3) hat nur dann eine Lösung, wenn $k=1$ ist, nämlich je
nachdem $D$ gerade oder ungerade, ist $x\equiv 0$ oder $\equiv q \pmod{2q}$,
dagegen keine Lösung, wenn $k>1$ ist, und dies gilt auch für
$q=2$. Denn in diesem Falle müßte $Q$ ungerade, $\Delta$ durch 4
teilbar sein, folglich $x$ gerade. Die Kongruenz

\[ \left(\frac{x}{2}\right)^2 \equiv Q^2\,\frac{\Delta}{4} \pmod{2^k} \]

ist aber nur lösbar für $k=1$, da $\Delta:4$ keine Diskriminante ist.

2) Ist $(D,q)=-1$, so ist die Kongruenz (3) nicht lösbar,
weil dann $D$ quadratischer Nichtrest von $q$ oder (für $q=2$)
von 8 ist.

3) Ist $(D,q)=+1$, so ist $D$ quadratischer Rest von $q$
(oder von 8) und die Kongruenz,

\[ x^2 \equiv D \mod 4q^k \]

hat, wie aus der Zahlentheorie bekannt ist, \gesperrt{zwei Wurzeln}.

Wir fassen diese Resultate übersichtlich so zusammen:

\[
\begin{array}{rcll}
\text{Ist }1) & (D,q)=0, & \text{so ist} & \psi(D,q)=1,\quad\psi(D,q^k)=0\quad(k>1),\\[0.4em]
\begin{array}{r}\text{\glqq\ }2)\\[0.3em]\text{\glqq\ }3)\end{array}
& \begin{array}{l}(D,q)=-1,\\[0.3em](D,q)=+1,\end{array}
& \begin{array}{c}\text{\glqq\ \glqq}\\[0.3em]\text{\glqq\ \glqq}\end{array}
& \left.\begin{array}{l}\psi(D,q^k)=0\\[0.3em]\psi(D,q^k)=2\end{array}\right\}\quad(k\geqq 1).
\end{array}
\]
% Source scan confirms the doubled greater-than-or-equal relation, transcribed as \geqq.

wozu noch kommt:

\[ \psi(D,1)=1. \]

Nach (5) ist also $\psi(D,m)$ immer dann gleich Null, wenn in
$m$ eine Primzahl aufgeht, für die $(D,q)=-1$ ist, oder wenn
in $m$ ein Primfaktor von $D$ mehr als einmal aufgeht.

In den anderen Fällen ist $\psi(D,q)$ eine Potenz von 2 % [sic] Gedruckt ist q; dem Zusammenhang nach wäre m zu erwarten., deren
Exponent gleich der Anzahl der in $m$, aber nicht in $D$ aufgehenden
 Primzahlen ist.

Zwischen den Symbolen $\psi(D,m)$ und $(D,m)$ besteht die
folgende allgemeine Beziehung immer unter der Voraussetzung,
daß $m$ relativ prim zu $Q$ ist:
%CONT>>
